%------------------------------------------------------------------------------%
\begin{question}
  Encontre a inversa da matriz
  \quad
  \(
    A =
    \begin{pmatrix}
      1 & 1 & 1 \\
      0 & 1 & 1 \\
      0 & 0 & 1
    \end{pmatrix}
  \)
\end{question}

%------------------------------------------------------------------------------%
\begin{resolution}{Matriz Aumentada}
  Matriz aumentada
  \[
    \begin{amatrix}[3]{3}
      1 & 1 & 1 & 1 & 0 & 0 \\
      0 & 1 & 1 & 0 & 1 & 0 \\
      0 & 0 & 1 & 0 & 0 & 1
    \end{amatrix}
  \]
\end{resolution}

%------------------------------------------------------------------------------%
\begin{resolution}{Escalonamento}
\begin{columns}
\column{0.3\linewidth}
  \[\;\;\;
    \begin{amatrix}[3]{3}
      1 & 1 & 1 & 1 & 0 & 0 \\
      0 & 1 & 1 & 0 & 1 & 0 \\
      0 & 0 & 1 & 0 & 0 & 1
    \end{amatrix}
  \]
  \pause
\column{0.7\linewidth}
  Eliminamos os elementos acima do terceiro pivô
  \\ \pause
  \(
    L_2 \leftarrow L_2 - L_3
  \)
  \\ \pause
  \(
    L_1 \leftarrow L_1 - L_3
  \)
\end{columns}
\vspace{-\baselineskip}
\[
  \hspace{-3ex}
  \pause
  L'_2
  \;=\;
  L_2 - L_3
  \pause
  \;=\;
  \begin{amatrix}[3]{3}
    0 & 1 & 1 & 0 & 1 & 0
  \end{amatrix}
  -
  \begin{amatrix}[3]{3}
    0 & 0 & 1 & 0 & 0 & 1
  \end{amatrix}
  \pause
  \;=\;
  \begin{amatrix}[3]{3}
    0 & 1 & 0 & 0 & 1 & -1
  \end{amatrix}
\]
\[
  \hspace{-3ex}
  \pause
  L'_1
  \;=\;
  L_1 - L_3
  \pause
  \;=\;
  \begin{amatrix}[3]{3}
    1 & 1 & 1 & 1 & 0 & 0
  \end{amatrix}
  -
  \begin{amatrix}[3]{3}
    0 & 0 & 1 & 0 & 0 & 1
  \end{amatrix}
  \pause
  \;=\;
  \begin{amatrix}[3]{3}
    1 & 1 & 0 & 1 & 0 & -1
  \end{amatrix}
\]
\pause
\[
  \begin{amatrix}[3]{3}
    1 & 1 & 0 & 1 & 0 & -1 \\
    0 & 1 & 0 & 0 & 1 & -1 \\
    0 & 0 & 1 & 0 & 0 & 1
  \end{amatrix}
\]
\end{resolution}

%------------------------------------------------------------------------------%
\begin{resolution}{Escalonamento}
\begin{columns}
\column{0.3\linewidth}
  \[\;\;\;
    \begin{amatrix}[3]{3}
      1 & 1 & 0 & 1 & 0 & -1 \\
      0 & 1 & 0 & 0 & 1 & -1 \\
      0 & 0 & 1 & 0 & 0 & 1
    \end{amatrix}
  \]
  \pause
\column{0.7\linewidth}
  Eliminamos os elementos acima do terceiro pivô
  \\ \pause
  \(
    L_1 \leftarrow L_1 - L_2
  \)
\end{columns}
\vspace{-\baselineskip}
\[
  \pause
  L'_1
  \;=\;
  L_1 - L_2
  \pause
  \;=\;
  \begin{amatrix}[3]{3}
    1 & 1 & 0 & 1 & 0 & -1
  \end{amatrix}
  -
  \begin{amatrix}[3]{3}
    0 & 1 & 0 & 0 & 1 & -1
  \end{amatrix}
  \pause
  \;=\;
  \begin{amatrix}[3]{3}
    1 & 0 & 0 & 1 & -1 & 0
  \end{amatrix}
\]
\pause
\[
  \begin{amatrix}[3]{3}
    1 & 0 & 0 & 1 & -1 & 0  \\
    0 & 1 & 0 & 0 & 1  & -1 \\
    0 & 0 & 1 & 0 & 0  & 1
  \end{amatrix}
\]
\end{resolution}

%------------------------------------------------------------------------------%
\begin{resolution}{Solução}
  \[
    \begin{amatrix}[3]{3}
      1 & 0 & 0 & 1 & -1 & 0  \\
      0 & 1 & 0 & 0 & 1  & -1 \\
      0 & 0 & 1 & 0 & 0  & 1
    \end{amatrix}
  \]

  \pause
  Como a matriz tem posto completo

  \pause
  A inversa existe e pode ser lida do lado direito
  \pause
  \[
    A^{-1} =
    \begin{bmatrix}
      1 & -1 & 0  \\
      0 & 1  & -1 \\
      0 & 0  & 1
    \end{bmatrix}
  \]
\end{resolution}

%------------------------------------------------------------------------------%
\endinput
